% Chapter — Composed environment force model (Phase 3 integration layer).
% Follows the mandatory FR-29 six-section template: purpose; assumptions
% and domain of validity; derivation (here: composition and plumbing, with
% the physics cited to the per-model chapters); discretization and
% implementation; validation evidence; references. The equation label
% eq:env:sum is echoed in cpp/src/models/environment.cpp.
\chapter{Composed Environment Force Model}
\label{ch:environment}

\section{Purpose}
\label{sec:env:purpose}

This chapter documents the Phase~3 force-composition layer
(\texttt{cpp/src/models/environment.cpp}): the point-mass translational
equations of motion about one central body (Earth, Moon, or Mars) under
the FR-5..FR-9 environment models. No physics is derived here --- every
term's derivation, domain, and golden-vector evidence live in the
per-model chapters (gravity, Chapter~\ref{ch:gravity}; third body,
Chapter~\ref{ch:thirdbody}; SRP and shadow, Chapter~\ref{ch:srp};
atmospheres, Chapters~\ref{ch:ussa76}, \ref{ch:harrispriester},
and~\ref{ch:marsatmosphere}; drag, Chapter~\ref{ch:drag}). What is
normative here, and nowhere else, is the \emph{composition}: the fixed
summation order, the frame and time plumbing each term is evaluated
through, the ephemeris body composition, the constants each term is wired
to, and the documented approximations those choices introduce.

\section{Assumptions and domain of validity}
\label{sec:env:domain}

\begin{enumerate}
  \item \textbf{Point-mass spacecraft.} The state is $(\vv{r}, \vv{v})$
    in the GCRF-oriented frame centered at the configured central body;
    attitude dynamics land in Phase~4. SRP and drag use the mass-normalized
    cannonball products $C_R A / m$ and $C_D A / m$ from the minimal
    \texttt{[spacecraft]} block.
  \item \textbf{One central body per run.} SOI-event origin re-centering
    (FR-3) is not part of Phase~3; a mission is Earth-, Moon-, or
    Mars-centered for its whole duration.
  \item \textbf{Inherited frame bounds.} Earth-fixed rotations use the
    IAU 2006/2000B CIO chain with polar motion neglected and constant
    $\Delta$UT1 $= 0$ (bounds in Chapter~\ref{ch:frames}); the Moon
    principal-axis frame interpolates the DE440 libration angles; the
    Mars frame uses the IAU~2015 rotational
    elements~\cite{archinal2018}.
  \item \textbf{Mars system barycenter.} DE440 stores the Mars
    \emph{system} barycenter. Using it as the Mars body center offsets
    the origin by the Phobos/Deimos barycenter displacement, bounded by
    the mass ratios ($\lesssim 2\times10^{-4}$~\unit{\metre}); the Mars
    system GM (GM4) is used consistently for both the central point-mass
    tier and the Mars third body.
  \item \textbf{Documented composition approximations} (each marked at
    its site below): the multi-occulter illumination product, the USSA76
    ceiling rule, and the Mars drag altitude datum.
\end{enumerate}

\section{Composition}
\label{sec:env:composition}

The total acceleration is the fixed-order sum (D-10: fixed
force-summation order, single-threaded)
\begin{equation}
  \label{eq:env:sum}
  \vv{a} \;=\; \vv{a}_{\mathrm{grav}}
  \;+\; \sum_{b\,\in\,\mathcal{B}} \vv{a}_{3,b}
  \;+\; \vv{a}_{\mathrm{SRP}}
  \;+\; \vv{a}_{\mathrm{drag}},
\end{equation}
with the third-body set $\mathcal{B}$ summed in the canonical order Sun,
Earth, Moon, Venus, Mars, Jupiter (the enabled subset, never the central
body). The Python validator records the enabled set in this same order in
the resolved configuration, so the summation order is visible in the
reproducibility anchor.

\paragraph{(a) Central-body gravity.} The \texttt{pointmass} tier applies
$\vv{a} = -\mu\,\vv{r}/r^3$ (Chapter~\ref{ch:twobody}) with $\mu$ from the
single constants home: the IERS value for Earth~\cite{iers2010} and the
DE440 header values for Moon and Mars~\cite{park2021}. The \texttt{j2} and
\texttt{harmonic} tiers rotate the position into the body-fixed frame,
evaluate the Pines model of Chapter~\ref{ch:gravity} on the loaded SRGRAV
field (whose own self-consistent $GM$/$R$/coefficient triple is used, per
that chapter), and rotate the acceleration back:
$\vv{a} = \mat{C}^{\mathsf{T}}\,
\vv{a}_{\mathrm{bf}}(\mat{C}\,\vv{r})$ with
$\mat{C} = \dcm{\mathrm{GCRF}}{\mathrm{bf}}$ from Chapter~\ref{ch:frames}.

\paragraph{(b) Third bodies.} Each enabled perturber contributes the
Battin $f(q)$ differential acceleration of Chapter~\ref{ch:thirdbody},
with the perturber position relative to the central body composed from
the trimmed DE440 ephemeris (Chapter~\ref{ch:ephemeris}) by exact segment
arithmetic --- never refit: SSB-relative bodies difference through the
central body's SSB position (e.g.\ the geocentric Sun is
$\vv{r}_{\mathrm{sun}} - (\vv{r}_{\mathrm{emb}} +
\vv{r}_{\mathrm{earth}})$); the Earth--Moon pair differences the
EMB-relative segments directly (the validated
\texttt{moon\_geocentric} composition). GM values are the DE440 header
constants (Table~\ref{tab:env:gm}), so each acceleration is consistent
with the DE440 positions it is evaluated at.

\begin{table}[htbp]
  \centering
  \caption{Third-body gravitational parameters: DE440 header values
    (constants block of the pinned \texttt{de440s.bsp}; published in
    Park et al.~\cite{park2021}), the single home
    \texttt{star/constants.hpp}. Mars and Jupiter are system GMs,
    matching the barycenter segments.}
  \label{tab:env:gm}
  \begin{tabular}{ll}
    \hline
    Body & $GM$ [\unit{\metre\cubed\per\second\squared}] \\
    \hline
    Sun & $1.32712440041279419\times10^{20}$ \\
    Earth & $3.98600435507\times10^{14}$ \\
    Moon & $4.902800118\times10^{12}$ \\
    Venus & $3.24858592\times10^{14}$ \\
    Mars system & $4.2828375816\times10^{13}$ \\
    Jupiter system & $1.267127641\times10^{17}$ \\
    \hline
  \end{tabular}
\end{table}

The DE440 Earth GM ($3.98600435507\times10^{14}$) differs from the IERS
TN36 central-body value ($3.986004418\times10^{14}$) by 7 parts in
$10^{9}$. This split is deliberate: the IERS value is the Phase~1 home
for Earth-centered two-body dynamics (frozen by the determinism record),
while third-body accelerations use the ephemeris's own constants. At
third-body magnitudes the difference is physically negligible
($<10^{-14}$~\unit{\metre\per\second\squared} at LEO).

\paragraph{(c) SRP.} The Sun position comes from the same ephemeris
composition; the illumination fraction over the configured occulters
(central body at minimum) combines per-occulter fractions of
Chapter~\ref{ch:srp} as a product,
$\nu = \prod_i \nu_i$ --- exact for a single occulter and whenever any
$\nu_i \in \{0, 1\}$, and a documented approximation when two bodies
partially occult the solar disk simultaneously (it assumes independent
disk coverage; the true value lies in
$[\max(0, \nu_1{+}\nu_2{-}1),\, \min(\nu_1,\nu_2)]$). Occulting-disk
radii: WGS84 equatorial radius for Earth~\cite{nima2000}, the IAU 2015
mean radius \qty{1737.4}{\kilo\metre} for the Moon and equatorial radius
\qty{3396.19}{\kilo\metre} for Mars~\cite{archinal2018}.

\paragraph{(d) Drag.} The air-relative velocity follows the FR-8
co-rotating rule (eq:drag:vrel, Chapter~\ref{ch:drag}):
$\vv{v}_{\mathrm{rel}} = \vv{v} - \vv{\omega} \times \vv{r}$ with
$\vv{\omega} = \omega\,\hat{\vv{z}}_{\mathrm{bf}}$ resolved in GCRF
($\hat{\vv{z}}_{\mathrm{bf}}$ is the third row of
$\dcm{\mathrm{GCRF}}{\mathrm{bf}}$; $\omega$ from
\texttt{star/constants.hpp}: IERS for Earth~\cite{iers2010}, the IAU 2015
spin rate for Mars~\cite{archinal2018}). Density selection:
\begin{itemize}
  \item \textbf{USSA76} (Earth): geodetic altitude over the WGS84
    ellipsoid~\cite{nima2000} as the geometric-altitude argument of
    Chapter~\ref{ch:ussa76}. Above the model's \qty{1000}{\kilo\metre}
    ceiling the composition takes $\rho = 0$ exactly (documented
    approximation mirroring the Harris--Priester ceiling, so drag
    vanishes instead of aborting; the omitted density is below
    $10^{-15}$~\unit{\kilogram\per\metre\cubed}).
  \item \textbf{Harris--Priester} (Earth): geodetic altitude over WGS84;
    the diurnal-bulge apex of Chapter~\ref{ch:harrispriester} built from
    the ephemeris geocentric Sun direction, with the $+30^{\circ}$
    rotation applied about the GCRF $+z$ axis and $\cos\psi$ formed with
    the GCRF position direction (the frame convention of the Orekit
    reference configuration; the GCRF-vs-true-equator axis difference is
    $\sim 20$~\unit{\mas}-scale and far inside the model's own
    uncertainty).
  \item \textbf{Mars exponential}: geodetic altitude over the IAU 2015
    Mars ellipsoid ($a = \qty{3396.19}{\kilo\metre}$,
    $1/f = 169.894$~\cite{archinal2018}) as the altitude argument of
    Chapter~\ref{ch:marsatmosphere}. The source model states no altitude
    datum (PRD A-3, confidence low); the ellipsoidal datum is this
    project's documented choice.
\end{itemize}

\section{Discretization and implementation}
\label{sec:env:implementation}

Time flows from the mission epoch through the Phase~2 time systems
(Chapter~\ref{ch:time}): the Python layer converts the UTC epoch to a
two-part TAI epoch across the binding (the core never parses text, D-2);
each right-hand-side evaluation at mission-elapsed $t$ forms
$\mathrm{TAI} + t$, from which the frame chain takes TT and the ephemeris
lookups take TDB seconds since J2000. The TDB epoch and the body-fixed
rotation are computed once per evaluation and shared by every term.

The model is evaluated through either integrator of
Chapter~\ref{ch:integrators} (FR-11): fixed-step RK4 with truth records
by step decimation (exactly the Phase~1 semantics), or adaptive RKF7(8)
through the event-aware driver of Chapter~\ref{ch:events} with truth
records sampled from the cubic Hermite dense output at
$k/f_{\mathrm{truth}}$. Dense-output interpolation error is bounded by
$h^{4}/384\,\max\lVert\vv{r}^{(4)}\rVert$ (Chapter~\ref{ch:integrators}),
so $h_{\max}$ is the logging-accuracy control; the cross-tool missions
pin $h_{\max} = \qty{30}{\second}$ (interpolation
$\sim\qty{2e-2}{\metre}$ at LEO, far below the exit-criterion gates).
Construction loads the gravity field and ephemeris before the time loop;
evaluation performs no heap allocation and no file I/O (D-10).

\section{Validation evidence}
\label{sec:env:validation}

\begin{table}[htbp]
  \centering
  \caption{Validation evidence for the composition layer (doctest,
    \texttt{cpp/tests/test\_environment.cpp}; pytest,
    \texttt{tests/python/}; verify checks,
    \texttt{python/star\_reacher/verify.py}). The physics of each term is
    gated in its own chapter's golden suite.}
  \label{tab:env:validation}
  \begin{tabular}{lp{8.2cm}}
    \hline
    Test & What it pins \\
    \hline
    \testid{env_composition_matches_sum_of_terms} &
      Full Earth stack (harmonic gravity, Sun+Moon third bodies, SRP,
      Harris--Priester drag) equals the independently recomposed sum of
      the per-model terms with the same frame chain and ephemeris
      arithmetic, to $10^{-12}$ relative. \\
    \testid{env_vrel_corotating_definition} &
      The isolated drag term uses
      $\vv{v} - \vv{\omega}\times\vv{r}$ (and demonstrably not the
      inertial velocity). \\
    \testid{env_run_double_run_bit_identity} &
      \texttt{run\_env} double runs are byte-identical for both
      integrators; RK4 and RKF7(8) final states agree at the metre level
      on a short arc. \\
    \testid{env_error_paths} &
      Constructor and configuration rejections: unknown bodies, central
      body as perturber, missing ephemeris, regime-inconsistent
      atmospheres, malformed adaptive controls; \texttt{gm()} wiring to
      the cited constants. \\
    \testid{V018} &
      Verify-tier perturbed double-run SHA-256 bit-identity on
      synthesized field and ephemeris fixtures (both integrators). \\
    \testid{test_leo_gravity_8x8_runs_and_is_deterministic} &
      Committed 7-day cross-tool gravity mission: bit-identical double
      run, bounded orbit, drift-free J2-inclusive energy. \\
    \testid{test_leo_drag_hp_runs_and_loses_energy} &
      Committed 7-day cross-tool drag mission: bit-identical double run,
      bounded orbit, secular energy loss from drag alone. \\
    \hline
  \end{tabular}
\end{table}

The external-truth comparison for the two committed cross-tool missions
(GMAT for gravity-only, Orekit for Harris--Priester drag; Phase~3 exit
criterion~5) is maintainer-side per D-15; the frozen configuration both
tools must replicate is specified in
\texttt{tests/golden/crosstool/README.md}.

\section{References}
\label{sec:env:references}

Model derivations: Chapters~\ref{ch:twobody}, \ref{ch:gravity},
\ref{ch:thirdbody}, \ref{ch:srp}, \ref{ch:ussa76},
\ref{ch:harrispriester}, \ref{ch:marsatmosphere}, and~\ref{ch:drag}.
Frames and time: Chapters~\ref{ch:frames} and~\ref{ch:time}. Integrators
and event-aware driver: Chapters~\ref{ch:integrators}
and~\ref{ch:events}. Constants: DE440 header GMs~\cite{park2021}; IERS
Conventions (2010)~\cite{iers2010}; IAU 2015 cartographic
report~\cite{archinal2018}; WGS84~\cite{nima2000}; formulation context in
Montenbruck and Gill~\cite{montenbruck2000} and
Vallado~\cite{vallado2013}.
